\chapter{Introducción}

\section{Contexto General y Planteamiento del Problema}

La electroencefalografía (EEG) ha sido, durante décadas, una de las herramientas fundamentales de la neurociencia clínica y experimental. Su alta resolución temporal —en el orden de los milisegundos—, su naturaleza no invasiva y su costo relativamente bajo frente a técnicas como la resonancia magnética funcional o la tomografía por emisión de positrones la han consolidado como el estándar para estudiar la dinámica cerebral en tiempo real \cite{buzsaki2006rhythms}. Sus aplicaciones abarcan el diagnóstico de epilepsia y trastornos del sueño, la investigación en neurociencia cognitiva y el desarrollo de interfaces cerebro-computador para comunicación asistida y neurorrehabilitación.

Ahora bien, la adquisición de EEG es un proceso frágil. En la práctica clínica y de investigación, la pérdida parcial o total de información en uno o más electrodos es frecuente. Estas fallas pueden deberse al aumento de la impedancia de contacto por transpiración o secado del gel conductor, a fallas en el amplificador como la saturación de voltaje, a movimientos bruscos del paciente que desconectan los sensores, o simplemente al descarte deliberado de canales que presentan una relación señal-ruido inaceptablemente baja durante la etapa de preprocesamiento \cite{nunez2006electric}.

La pérdida de canales no es un problema cosmético: compromete la integridad de todo análisis posterior. La ausencia de información espacial distorsiona la topografía de la señal, afecta la localización de fuentes neuronales y altera biomarcadores electrofisiológicos como los Potenciales Relacionados a Eventos (ERP) y la Densidad Espectral de Potencia (PSD)~\cite{luck2014introduction}. En sistemas modernos de aprendizaje automático aplicados a interfaces cerebro-computador, la pérdida impredecible de canales cambia la dimensionalidad de entrada de los clasificadores y obliga, con frecuencia, a descartar ensayos completos de datos clínicamente valiosos \cite{lotte2018review}. La interpolación precisa de estos canales es, por tanto, un paso ineludible del flujo de preprocesamiento, y su calidad condiciona directamente la validez de cualquier conclusión neurofisiológica o clínica posterior.

\section{Limitaciones de los Enfoques Tradicionales}

Históricamente, el problema de la interpolación de canales EEG ha sido abordado mediante métodos espaciales o estadísticos. El enfoque más extendido y utilizado por defecto en bibliotecas de análisis como EEGLAB \cite{delorme2004eeglab} o MNE-Python \cite{gramfort2013meg,gramfort2014mne} es la interpolación por \textit{spherical splines} \cite{perrin1989}. Este método, junto con otros enfoques basados en distancia euclidiana como la interpolación por vecinos más cercanos o las funciones de base radial, asume que la cabeza humana es una esfera conductora perfecta y estima el voltaje del canal faltante proyectando geométricamente los voltajes de los electrodos circundantes sobre dicha superficie.

Aunque estos métodos son rápidos y matemáticamente elegantes, adolecen de una limitación fundamental: operan estrictamente en el dominio espacial estático para cada instante de tiempo individual. Al hacerlo, ignoran dos realidades neurofisiológicas esenciales. En primer lugar, la distribución espacial del potencial eléctrico en el cuero cabelludo no es isotrópica ni homogénea: la conducción volumétrica a través de tejidos de distinta conductividad ---cuero cabelludo, cráneo, líquido cefalorraquídeo--- distorsiona el campo de manera no uniforme, de modo que un modelo esférico simplificado no captura completamente esta heterogeneidad \cite{nunez2006electric}. En segundo lugar, las señales EEG son generadas por osciladores neuronales sujetos a leyes físicas de continuidad y momento. Al reconstruir un canal muestra por muestra de forma aislada, los métodos puramente espaciales son incapaces de preservar la dinámica de alta frecuencia, generando trazos con ruido de fase y distorsiones en los picos transitorios \cite{lai2005controlled}.
Los métodos estadísticos de descomposición, Análisis de Componentes Independientes (ICA) y Análisis de Componentes Principales (PCA), proyectan las señales observadas sobre fuentes latentes y reconstruyen el canal perdido a partir de esa proyección \cite{jung2001imaging}. Sin embargo, cuando varios canales contiguos fallan simultáneamente, estos métodos tienden a inyectar ruido global o a perder los picos transitorios locales que definen la morfología fina de los ERP.

\section{El Advenimiento del Procesamiento de Señales en Grafos}

Para superar estas limitaciones, en los últimos años ha emergido con fuerza el marco del Procesamiento de Señales en Grafos (GSP, por sus siglas en inglés) \cite{shuman2013,ortega2018}. El GSP ofrece una perspectiva natural y poderosa para lidiar con datos que residen en dominios irregulares no euclidianos. En lugar de modelar la cabeza como un espacio continuo tridimensional, GSP modela la red de electrodos como un grafo discreto $\mathcal{G} = (\mathcal{V}, \mathcal{E}, \mathbf{W})$, donde cada electrodo es un nodo ($\mathcal{V}$), las aristas ($\mathcal{E}$) definen qué pares de electrodos están conectados, y la matriz de pesos ($\mathbf{W}$) cuantifica la intensidad de cada conexión. Las relaciones de similitud o distancia fisiológica entre electrodos se codifican en $\mathbf{W}$ mediante núcleos de distancia gaussianos o basados en correlación.
Bajo este paradigma, el voltaje medido en un instante se interpreta como una señal de grafo, y conceptos clásicos como la Transformada de Fourier, el filtrado paso-bajo y el muestreo se trasladan al dominio del grafo mediante el operador Laplaciano $\mathbf{L}$. La intuición es simple: las señales en nodos fuertemente conectados tienden a ser similares. Al minimizar una función de costo que penaliza las variaciones abruptas sobre la topología del grafo, los datos faltantes se reconstruyen con mayor fidelidad que usando solo la distancia euclidiana \cite{narang2013}.

Una extensión aún más avanzada es la regularización espacio-temporal. Trabajos recientes han introducido enfoques como la Regularización Temporal de Suavidad de Sobolev (TRSS, por sus siglas en inglés) \cite{giraldo2022}, que optimiza conjuntamente la topología espacial y la derivada temporal de la señal. Al penalizar saltos bruscos en el tiempo de forma simultánea a los saltos en el espacio, TRSS impone restricciones fisiológicas reales: las neuronas no cambian su potencial de membrana instantáneamente, sino que obedecen dinámicas de integración continua. Esta restricción dual —espacial y temporal— es lo que la teoría predice como el factor diferencial frente a los métodos puramente espaciales.

Aunque existen enfoques de aprendizaje profundo para imputación de series temporales y representación de señales, esta tesis no los incorpora como línea experimental. La razón es metodológica: el objetivo es comparar métodos algorítmicos sin entrenamiento —geométricos, estadísticos, GSP y regularización espacio--temporal— bajo un protocolo controlado y trazable. Incluir modelos entrenables requeriría particiones de entrenamiento, validación externa y análisis de generalización que exceden el alcance de este trabajo.

\section{La Brecha en la Literatura y la Motivación del Trabajo}

A pesar del potencial teórico de los algoritmos GSP y TRSS, su validación empírica en la literatura de señales biomédicas presenta limitaciones metodológicas relevantes. Se identifican tres brechas que motivan el diseño experimental de esta tesis.

El primer problema es que los estudios que proponen nuevos algoritmos GSP suelen compararse contra \textit{spherical splines} usando las configuraciones por defecto de los paquetes de software, sin realizar una búsqueda rigurosa de hiperparámetros para los métodos de referencia. Esto puede inflar artificialmente las ventajas del método propuesto y distorsionar la percepción del verdadero estado del arte \cite{saldivia2026benchmark}. El segundo problema es que la mayoría de los trabajos restringe su evaluación a métricas de error de amplitud global —como el RMSE o el MAE— y aplica escenarios de desconexión simples, tales como la remoción de canales aleatorios aislados en bases de datos pequeñas. Rara vez se investiga el impacto morfológico de la reconstrucción mediante métricas temporales como el \textit{Dynamic Time Warping} o espectrales como la Distancia Espectral Logarítmica. El tercer problema es la escasez de comparaciones contra implementaciones comunitarias maduras. En particular, la interpolación de MNE-Python incorpora más de una década de refinamientos de ingeniería —normalización de coordenadas, regularización diagonal adaptativa, aumentación del sistema con restricciones lineales, pseudoinversa con tolerancia adaptativa y un mecanismo de reintento automático— que la hacen sustancialmente más robusta que una implementación ingenua de \textit{spherical splines}. Ignorar esta implementación equivale a comparar contra un método de referencia debilitado artificialmente.

La motivación concreta de este trabajo nace directamente de estas brechas. El objetivo no es diseñar desde cero un nuevo marco teórico matemático, sino someter a la frontera del conocimiento actual ---métodos geométricos, GSP y regularización temporal--- a una evaluación empírica rigurosa, exhaustiva y justa. Al utilizar Optuna \cite{akiba2019optuna} para explorar configuraciones comparables en múltiples conjuntos de datos, esta tesis busca identificar el punto de cruce práctico entre el paradigma clásico y el paradigma de grafos: cuándo basta una referencia madura como MNE-Python y cuándo la regularización espacio--temporal aporta valor adicional.
\subsection{La pregunta central: ¿por qué funciona tan bien la interpolación de MNE-Python?}
\label{sec:why_mne_works}

Una pregunta recurrente que emerge del análisis preliminar —y que motiva una sección dedicada en el marco teórico de esta tesis— es por qué la función de interpolación de MNE-Python produce resultados tan competitivos, incluso frente a métodos teóricamente superiores como TRSS. La respuesta, que se desarrolla en profundidad en el Capítulo~2, reside en que la implementación de MNE no es una simple transcripción del artículo original de Perrin y colaboradores \cite{perrin1989}, sino que incorpora un conjunto de mejoras de ingeniería acumuladas durante más de diez años de desarrollo como código abierto.

Estas mejoras —normalización de coordenadas a esfera unitaria, uso global de canales, regularización diagonal adaptativa, aumentación del sistema con restricciones lineales, pseudoinversa con tolerancia adaptativa y reintento automático— no aparecen descritas en la literatura académica, pero están presentes en el código fuente de MNE. Son las responsables de que un método formulado en 1989 siga siendo hoy una referencia competitiva. Esta tesis dedica una sección completa del marco teórico (Capítulo~2) a documentar estos mecanismos, y una validación comparativa independiente (Capítulo~4) a cuantificar la brecha real entre TRSS y MNE.

\section{Objetivos}

\subsection{Objetivo General}

Comparar el rendimiento espacial, temporal y espectral de los métodos de reconstrucción de canales EEG basados en GSP con regularización temporal frente a las referencias geométricas bajo un marco experimental riguroso, reproducible y optimizado mediante optimización bayesiana. 
\subsection{Objetivos Específicos}

\begin{itemize}
  \item \textbf{OE1 — Sistema unificado de evaluación.} Implementar dieciocho métodos de interpolación en un sistema modular con interfaz común, definir protocolos de pérdida aleatoria y contigua (10\%--40\%), e integrar Optuna para la optimización imparcial de todos los algoritmos, incluyendo una fase de selección preliminar y reducción de candidatos antes de la comparación final.

  \item \textbf{OE2 — Evaluación multi-métrica y fisiológica.} Cuantificar el desempeño mediante seis métricas primarias (MAE, RMSE, SNR, DTW, LSD, Coherencia) sobre tres conjuntos de datos y 120 escenarios de pérdida, documentando tanto los resultados preliminares de selección método--grafo como la preservación morfológica y espectral mediante métricas auxiliares, visualizaciones temporales y potencia por banda.

  \item \textbf{OE3 — Trazabilidad y reproducibilidad.} Generar un repositorio de artefactos computacionales —código, configuraciones, registros de optimización y resultados— que garantice la trazabilidad completa de cada afirmación cuantitativa hasta el archivo que la origina.
\end{itemize}

\section{Preguntas de Investigación}

En lugar de formular la evaluación como una prueba de hipótesis estadística, esta tesis organiza el análisis en torno a dos preguntas comparativas:

\begin{quote}
\textbf{P1.} ¿En qué condiciones la regularización espacio--temporal mejora la reconstrucción de canales EEG faltantes frente a referencias geométricas y a MNE-Python?
\end{quote}

\begin{quote}
\textbf{P2.} ¿Qué compromiso práctico existe entre precisión de reconstrucción, fidelidad espectral y tiempo de cómputo al elegir entre TRSS y MNE-Python?
\end{quote}

Estas preguntas se responden mediante comparaciones pareadas, estratificación por escenario y análisis multi-métrico; no se interpreta la tesis como aceptación o rechazo de hipótesis nulas.
\section{Alcance y Aportes de la Investigación}

El alcance de este trabajo abarca la ejecución computacional a gran escala del sistema de evaluación sobre tres conjuntos de datos, ciento veinte escenarios de falla y dieciocho métodos de interpolación. A diferencia de las evaluaciones acotadas que predominan en la literatura, aquí se consolidan miles de ensayos EEG bajo más de quince mil configuraciones paramétricas exploradas mediante Optuna.

Para garantizar la reproducibilidad, la ejecución se ancla en un registro exhaustivo de corridas experimentales que fija explícitamente las combinaciones de conjuntos de datos, modos de pérdida, semillas, tipos de grafo y familias de métodos. Este diseño evita las evaluaciones improvisadas y permite trazar cada afirmación cuantitativa reportada en los capítulos posteriores hasta el artefacto experimental que la origina, según se detalla en el Apéndice~D.

Los aportes principales se articulan en cuatro ejes. Primero, se proporciona evidencia empírica pareada de que TRSS mejora métricas de amplitud y ajuste temporal frente a referencias geométricas en regímenes difíciles, una vez que todos los métodos han sido optimizados. Segundo, se muestra que bajo pérdida contigua severa (30\% o más de los canales) la restricción temporal puede compensar la pérdida de información espacial local. Tercero, se documentan en formato académico los mecanismos de ingeniería que convierten a MNE-Python en una referencia inesperadamente competitiva. Cuarto, se entrega un sistema unificado en Python, rastreable y reproducible para ejecutar interpolaciones GSP con optimización bayesiana automatizada sobre datos EEG.

\section{Estructura de la Tesis}

Para desarrollar el análisis propuesto, el documento se organiza en seis capítulos. El Capítulo~2 presenta los fundamentos neurofisiológicos del EEG, el Procesamiento de Señales en Grafos, los métodos de interpolación y los mecanismos internos de MNE-Python. El Capítulo~3 define la metodología: datos, máscaras de pérdida, construcción de grafos, métodos, métricas, arquitectura, trazabilidad y separación entre fase exploratoria y fase comparativa. El Capítulo~4 documenta los experimentos y resultados: selección preliminar, reducción de candidatos, optimización intermedia, comparación final TRSS--MNE, señales representativas, PSD, contribución temporal y frontera práctica de uso. El Capítulo~5 discute las causas y limitaciones de los resultados. El Capítulo~6 resume conclusiones y trabajo futuro.

Los apéndices complementan el documento con derivaciones matemáticas, tablas extensas, protocolo de reproducibilidad, inventario de artefactos y material histórico. Esta organización permite seguir una cadena de decisiones verificables: qué métodos y grafos se consideraron, cuáles se descartaron, qué variantes pasaron a la comparación final y cómo se interpretan los resultados sin depender de una prueba de hipótesis como criterio único.

